\begin{proofbox}
	\textbf{\textcolor{CProof}{思路提示}}

	利用定义：平行线间距离等于一条直线上任意一点到另一条直线的距离。因此，在 $l_1$ 上任取一点 $P$，求 $P$ 到 $l_2$ 的距离，化简即得与点无关的表达式。

	\medskip
	\textbf{\textcolor{CProof}{正式推导}}
	\begin{description}[style=nextline, leftmargin=2.8em, labelsep=0.5em, itemsep=0.55em]
		\item[\textbf{步骤一（取点代入）：}] 设 $P(x_0,y_0)$ 满足 $Ax_0+By_0+C_1=0$。
			\par\textbf{理由：} 由定义，距离与点选取无关，可任取一点。

		\item[\textbf{步骤二（距离公式）：}] 点 $P$ 到 $l_2$ 的距离为
			\[
				d = \frac{|Ax_0+By_0+C_2|}{\sqrt{A^{2}+B^{2}}}.
			\]
			\par\textbf{理由：} 点到直线距离公式 $d = \dfrac{|Ax_0+By_0+C|}{\sqrt{A^{2}+B^{2}}}$。

		\item[\textbf{步骤三（代入化简）：}] 由 $Ax_0+By_0 = -C_1$，代入上式得
			\[
				\begin{aligned}
					d &= \frac{| -C_1 + C_2 |}{\sqrt{A^{2}+B^{2}}} \\
					&= \frac{|C_1-C_2|}{\sqrt{A^{2}+B^{2}}}.
			\end{aligned}
			\]
			\par\textbf{理由：} 代入后分子与 $P$ 无关，只取决于常数项。

		\item[\textbf{步骤四（得出公式）：}] 上式即为两平行线间距离公式。
			\par\textbf{理由：} 表达式不含 $x_0,y_0$，公式成立。
	\end{description}

	\medskip
	\textbf{\textcolor{CProof}{结论回扣}}

	因此，两平行直线 $l_1:Ax+By+C_1=0$ 与 $l_2:Ax+By+C_2=0$（$C_1\neq C_2$）间的距离为 $d = \dfrac{|C_1-C_2|}{\sqrt{A^{2}+B^{2}}}$。
\end{proofbox}
